% Part: first-order-logic
% Chapter: natural-deduction
% Section: quantifier-rules

\documentclass[../../../include/open-logic-section]{subfiles}

\begin{document}

\olfileid[tr]{fol}{ntd}{qrl}

\olsection{Niceleyici Kuralları}

\subsection{$\lforall$ Kuralları}

\begin{defish}
\AxiomC{$!A(a)$}
\RightLabel{\Intro{\lforall}}
\UnaryInfC{$\lforall[x][\Atom{!A}{x}]$}
\DisplayProof
\hfill
\AxiomC{$\lforall[x][\Atom{!A}{x}]$}
\RightLabel{\Elim{\lforall}}
\UnaryInfC{$!A(t)$}
\DisplayProof
\end{defish}

$\lforall$ kurallarında $t$, kapalı bir terimdir (hiçbir değişken içermeyen
bir terim); $a$ ise !!a{constant} olarak seçilir. Bu sabit,
sonuç~$\lforall[x][!A(x)]$ içinde ya da öncül~$!A(a)$ bakımından
!!{undischarged} kalıp bu öncülle biten !!{derivation} içinde yer alan
herhangi bir varsayımda geçmez. $a$'ya \Intro{\lforall}
çıkarımının \emph{özdeğişkeni} deriz.\footnote{Yukarıdaki kuralda $a$ bir
sabit olduğu hâlde ``özdeğişken'' terimini kullanıyoruz. Bunun tarihsel
nedenleri vardır.}

\subsection{$\lexists$ Kuralları}

\begin{defish}
\AxiomC{$\Atom{!A}{t}$}
\RightLabel{\Intro{\lexists}}
\UnaryInfC{$\lexists[x][\Atom{!A}{x}]$}
\DisplayProof
\hfill
\AxiomC{$\lexists[x][\Atom{!A}{x}]$}
\AxiomC{[$\Atom{!A}{a}$]$^n$}
\DeduceC{$!C$}
\DischargeRule{\Elim{\lexists}}{n}
\BinaryInfC{$!C$}
\DisplayProof
\end{defish}

Burada da $t$ kapalı bir terimdir; $a$ ise !!a{constant} olarak seçilir. Bu
sabit öncül $\lexists[x][!A(x)]$ içinde, sonuç~$!C$ içinde veya
!!{undischarged} kalıp iki öncülle biten !!{derivation}s içinde yer alan
herhangi bir varsayımda ($!A(a)$ varsayımları dışında) geçmez. $a$'ya
\Elim{\lexists} çıkarımının
\emph{özdeğişkeni} deriz.

Bir özdeğişkenin, \Intro{\lforall} veya \Elim{\lexists} çıkarımında yukarıda
belirtilen yasaklı sonuç ya da ana öncül formüllerinde geçmemesi; ayrıca
!!{undischarged} kalan varsayımlarla öncüllere götüren !!{derivation}s içinde
yer almaması koşuluna
\emph{özdeğişken koşulu} denir.

\begin{explain}
$!A$ bir !!a{formula} olduğunda ve !!{variable}~$x$ bu formülde serbest
geçtiğinde, bunu
$!A(x)$ yazarak belirttiğimizi hatırlayın. Aynı bağlamda $!A(t)$,
$\Subst{!A}{t}{x}$ ifadesinin kısaltmasıdır. Dolayısıyla
$\Intro\lexists$ kuralını şöyle de yazabilirdik:
\begin{prooftree}
\AxiomC{$\Subst{!A}{t}{x}$}
\RightLabel{\Intro{\lexists}}
\UnaryInfC{$\lexists[x][!A]$}
\end{prooftree}
$t$'nin $!A$ içinde zaten geçebileceğine dikkat edin; örneğin $!A$,
$\Atom{\Obj P}{t,x}$ olabilir. Bu nedenle $\Atom{\Obj P}{t,t}$'den
$\lexists[x][\Atom{\Obj P}{t,x}]$ çıkarmak, $\Intro\lexists$ kuralının
doğru bir uygulamasıdır: öncüldeki $t$ geçişlerinin birini veya birkaçını,
ama zorunlu olarak hepsini değil, bağlı !!{variable}~$x$ ile
``değiştirebilirsiniz''. Ne var ki $\Intro\lforall$ ve~$\Elim\lexists$
kurallarındaki özdeğişken koşulları, !!{constant}~$a$'nın $!A$ içinde
geçmemesini gerektirir. Bu yüzden $\Atom{\Obj P}{a,a}$'dan
$\Intro\lforall$ kullanarak $\lforall[x][\Atom{\Obj P}{a,x}]$ sonucunu
doğru biçimde çıkaramazsınız.
\end{explain}

\begin{explain}
\Intro{\lexists} ve \Elim{\lforall} kurallarında başka hiçbir kısıtlama yoktur;
$t$ terimi herhangi bir kapalı terim olabilir, dolayısıyla ek bir koşulu gözetmemiz
gerekmez. Buna karşılık \Elim{\lexists} ve \Intro{\lforall} kurallarında
özdeğişken koşulu, !!{constant}~$a$'nın sonuçta veya herhangi bir
!!{undischarged} varsayımda hiçbir yerde geçmemesini gerektirir. Bu koşul,
sistemin sağlam olmasını güvence altına alır: sistemin !!{derive}s
eylemleriyle ulaştığı !!{sentence}s yalnızca mantıksal olarak çıktıkları
!!{undischarged} varsayımlara dayanabilir. Bu koşul
olmasaydı aşağıdaki çıkarıma izin
verilirdi:
\begin{prooftree}
  \AxiomC{$\lexists[x][!A(x)]$}
  \AxiomC{$\Discharge{!A(a)}{1}$}
  \RightLabel{*\Intro{\lforall}}
  \UnaryInfC{$\lforall[x][!A(x)]$}
  \RightLabel{\Elim{\lexists}}
  \BinaryInfC{$\lforall[x][!A(x)]$}
\end{prooftree}
Oysa $\lexists[x][!A(x)] \Entails/ \lforall[x][!A(x)]$.

Niceleyicilerin giderme kuralları, !!{variable}s yerine yalnızca kapalı
terimlerin konmasına izin verdiği için türetilebilen her !!{formula}, bir
!!{sentence}s kümesinden çıkıyorsa kendisi de !!a{sentence} olur.
\end{explain}


\end{document}
